\begin{abstract}
In cooperative multi-agent reinforcement learning with heterogeneous roles, agents must not only follow role-specific reward signals but also route attention to the right teammates at each timestep. We ask whether policy architecture determines how effectively role labels are exploited as routing signals. We propose a Label-Conditioned Intra-Set Self-Attention (SLIC) architecture in which role one-hots appear in both the agent's own token and each ally token, enabling role-pair-conditioned Query/Key routing in cross-attention. We evaluate five ablation conditions over 2M training steps using MAPPO on a custom 3-role MiniGrid environment (Strike, Vanguard, Support). The full architecture achieves 0.920 $\pm$ 0.054 win rate; removing role labels with a shared policy degrades performance by 8.5 percentage points (0.780 $\pm$ 0.113). Attention produces $2\times$ more structured role-routing patterns than MLP (routing matrix std: 0.323 vs.\ 0.153), while MLP underutilizes role labels by 2.4--3.7$\times$ relative to spatial features. Three of five a priori routing predictions are confirmed; two are revised but behaviorally interpretable. The contribution is an analysis framework\textemdash role-routing matrix, $\Delta_{\max}$ formation sensitivity, and gradient attribution comparison\textemdash that makes role-to-role routing legible across architectures.
\end{abstract}
